\begin{lemma}[Stability of the quadratic envelopes]
Let $F$ be a finite-dimensional real inner-product space, let $I\subseteq\R$ be
open and let $\Omega=\{(x,t):t\in I\}$. Let $f$ and $f_1,f_2,\ldots$ be real
functions on space-time, continuous on $\Omega$, with $f_j\to f$ uniformly on
every compact subset of $\Omega$, with $f$ of subquadratic spatial growth on
compact time intervals of $I$ in the sense of \noderef{SubquadraticGrowth} and
with the sequence having the common growth control of
\noderef{UnifSubquadraticGrowth} on $I$:
\[
  \lim_{M\to\infty}\sup_j\sup_{\substack{|y|\ge M\\t\in J}}
      \frac{|f_j(y,t)|}{|y|^2}=0\qquad(J\Subset I).
\]
Let $\ell>0$. Then $\Penv_\ell f_j\to\Penv_\ell f$ uniformly on every compact
set of centers $z=(x,t)$ with $t\in I$ and $t+\ell\in I$, and
$\Menv_\ell f_j\to\Menv_\ell f$ uniformly on every compact set of centers with
$t\in I$ and $t-\ell\in I$, the envelopes being those of \noderef{QuadEnvP} and
\noderef{QuadEnvM}.
\end{lemma}
\begin{proof}
Write $\Omega=\{(y,\sigma):\sigma\in I\}$, an open set on which $f$ and every
$f_j$ are continuous.

\emph{Step 1: uniform localization.}
Let $\mathcal K$ be a compact set of centers $z=(x,t)$ with $t\in I$ and
$t+\ell\in I$ for all $z\in\mathcal K$ (the case of the backward envelope is
identical with $t-\ell$ in place of $t+\ell$, and we indicate the changes at the
end). Put
\[
  R_0=\max\{|x|:(x,t)\in\mathcal K\},
  \qquad
  J=\{t:(x,t)\in\mathcal K\}\cup\{t+\ell:(x,t)\in\mathcal K\},
\]
so $R_0<\infty$ and $J$ is a compact subset of $I$, being the union of two
continuous images of $\mathcal K$.

Apply \noderef{UnifSubquadraticGrowth} with this $J$ and $\varepsilon=1/(8\ell)$:
there is $M_1>0$ with
\begin{equation}\label{eq:qes-growthj}
  |f_j(y,\sigma)|\le\frac{|y|^2}{8\ell}
  \qquad(j\in\mathbb N,\ |y|\ge M_1,\ \sigma\in J),
\end{equation}
and apply \noderef{SubquadraticGrowth} with the same $J$ and $\varepsilon$ to get
$M_2>0$ with the same bound for $f$. Put
$M_3=\max\{M_1,M_2,4R_0,1\}$.

For $|y|>M_3$, any $x$ with $|x|\le R_0$ and any $\sigma\in J$ we have
$|y-x|\ge|y|-R_0\ge\tfrac34|y|$, so, exactly as in the corresponding
computation for a single function,
\begin{equation}\label{eq:qes-far}
  g(y,\sigma)-\frac{|y-x|^2}{2\ell}
   \le\frac{|y|^2}{8\ell}-\frac{9|y|^2}{32\ell}=-\frac{5|y|^2}{32\ell}
\end{equation}
for $g=f$ and for every $g=f_j$.

The set $\mathcal C_3=\{(y,\sigma):|y|\le M_3,\ \sigma\in J\}$ is a compact
subset of $\Omega$. Since $f$ is continuous on $\Omega$, $N_f:=\max_{\mathcal
C_3}|f|$ is finite; and since $f_j\to f$ uniformly on $\mathcal C_3$, there is
$j_0$ with $\sup_{\mathcal C_3}|f_j-f|\le1$ for $j\ge j_0$, while for each of
the finitely many $j<j_0$ the continuous $f_j$ is bounded on $\mathcal C_3$.
Hence
\begin{equation}\label{eq:qes-N}
  N:=\sup_{j\in\mathbb N}\ \max_{\mathcal C_3}|f_j|\ \vee\ N_f<\infty .
\end{equation}
Put $M=\max\bigl\{M_3,\ \sqrt{\tfrac{32}{5}N\ell}+1\bigr\}$. For $|y|>M$,
$|x|\le R_0$ and $\sigma\in J$ we have $\tfrac{5|y|^2}{32\ell}>N$, so
\eqref{eq:qes-far} and \eqref{eq:qes-N} give, for $g=f$ and every $g=f_j$,
\[
  g(y,\sigma)-\frac{|y-x|^2}{2\ell}<-N\le g(x,\sigma),
\]
the last inequality because $|x|\le R_0\le M_3$. In particular no competitor
with $|y|>M$ exceeds the competitor $y=x$.

\emph{Step 2: both envelopes are maxima over one fixed compact ball.}
Let $z=(x,t)\in\mathcal K$. By Step 1 applied with $\sigma=t+\ell\in J$, the set
of competitor values in \noderef{QuadEnvP} for $g$ at $z$ and length $\ell$ is
nonempty and bounded above, and its supremum equals the maximum over
$\{|y|\le M\}$, which exists because $y\mapsto g(y,t+\ell)-|y-x|^2/(2\ell)$ is
continuous on that compact ball. Hence
\begin{equation}\label{eq:qes-max}
  \Penv_\ell g(z)=\max_{|y|\le M}
     \Bigl(g(y,t+\ell)-\frac{|y-x|^2}{2\ell}\Bigr)
  \qquad(g=f\text{ or }g=f_j,\ z\in\mathcal K).
\end{equation}

\emph{Step 3: conclusion.}
For bounded real functions $G,H$ on a nonempty set $Q$ one has
$|\sup_QG-\sup_QH|\le\sup_Q|G-H|$. Applying this to
\eqref{eq:qes-max} with $Q=\{|y|\le M\}$, the quadratic terms cancelling,
\[
  \bigl|\Penv_\ell f_j(z)-\Penv_\ell f(z)\bigr|
   \le\max_{|y|\le M}\bigl|f_j(y,t+\ell)-f(y,t+\ell)\bigr|
   \le\sup_{\mathcal C}\,|f_j-f| ,
\]
where $\mathcal C=\{(y,\sigma):|y|\le M,\ \sigma\in J\}$ is a compact subset of
$\Omega$, and the bound is independent of $z\in\mathcal K$. By hypothesis
$f_j\to f$ uniformly on $\mathcal C$, so the right side tends to $0$. Hence
$\Penv_\ell f_j\to\Penv_\ell f$ uniformly on $\mathcal K$.

For the backward envelope, repeat Steps 1--3 with $t-\ell$ in place of
$t+\ell$ in the definition of $J$, with $-f$ and $-f_j$ in place of $f$ and
$f_j$ in \eqref{eq:qes-far} — these satisfy the same hypotheses, the growth
conditions of \noderef{SubquadraticGrowth} and
\noderef{UnifSubquadraticGrowth} being invariant under a change of sign — and
with infima and minima in place of suprema and maxima, using
$\inf A=-\sup(-A)$ and \noderef{QuadEnvM}. This gives
$\Menv_\ell f_j\to\Menv_\ell f$ uniformly on every compact set of centers
$z=(x,t)$ with $t\in I$ and $t-\ell\in I$.
\end{proof}
